\documentclass{article}
\usepackage[a4paper,margin=1in]{geometry}
\usepackage{amsmath,amssymb}
\usepackage{mdframed}
\usepackage{xcolor}

\setlength{\parskip}{1em}  % Set the space between paragraphs
\setlength{\parindent}{0pt}  % Remove the indentation at the beginning of paragraphs
%lang=fr

\mdfsetup{
topline=false,
rightline=false,
bottomline=false,
linewidth=5pt,
innerleftmargin=15pt,
innerrightmargin=0pt,
skipabove=\baselineskip,
skipabove=1.2\baselineskip,
}

\definecolor{lightblue}{RGB}{160, 216, 239}
\definecolor{darkblue}{RGB}{40, 100, 158}

\begin{document}

\title{la musique et les mathematiques}
\author{B. Ruben}
\date{\today}

\maketitle

\section{Introduction}
La musique dépend des mathématiques, et les travaux de pythagore le montre. Une note de musique est définie par sa frequence fondamentale et ses harmoniques. La perception simultanée de plusieurs notes peut donner l'impression que les notes sonnent bien ensemble, on dit qu'elles sont consonantes

\subsection{TP}



\end{document}
